At a fixed state write
\[
 f_a(z)=V_1(F(x,a,z)),\qquad
 Q_x(a)=S(x)+h(a)+\beta\rho_\alpha(f_a).
\]
The actual complete-policy cost is $J_x(a)$, with the same implemented
last-date policy for every current action.  Its retained uniform guarantee is
\begin{equation}\label{eq:future120}
 J_x(a)=Q_x(a)+\eta_x(a),\qquad
 0\le\eta_x(a)\le\delta,\qquad
 \delta=\beta\varepsilon_1,\quad \varepsilon_1=2^{-20}.
\end{equation}
The function $\eta_x$ is not assumed convex or differentiable.

\subsection{A comparison modulus from the original technology}
The following constant concerns the original optimal continuation, not a
fitted critic or an empirical gradient bound.  This addition is restricted
to the original two-dimensional application.

\begin{theorem}[Original action comparison modulus]
\label{thm:modulus120}
For $d=2$, $\beta=15/16$, $R=1/32$, and
$\alpha\in\{1/2,3/4\}$, put $t=1-\alpha$ and
\begin{equation}\label{eq:modulus120}
 M_\alpha=\frac{1413}{128}
       +\beta\left\{\frac{127}{16}
                 +\frac{19625\alpha}{20480t}\right\}.
\end{equation}
The original $V_1$ has an $M_\alpha$-Lipschitz gradient on the state
cube, with the continuous boundary extension.  Consequently, for every
original state and every two feasible actions,
\begin{equation}\label{eq:mixedmodulus120}
 \operatorname{Lip}_z(f_b-f_a)
 \le E_{ab}:=\frac45M_\alpha|b-a|.
\end{equation}
Both conclusions include switches of the constrained terminal optimizer
and of its active risk region.
\end{theorem}

The proof normalizes terminal investment as $A(y)v$, $v\in[0,1]$,
and bounds the state semiconcavity of that fixed-domain objective.
It therefore does not differentiate a rounded actor or ignore the
capacity multiplier.  The earlier expected-cost smoothness result in
Section~\ref{sec:accuracy62} is retained with its original criterion;
the extra term in \eqref{eq:modulus120} pays for the movement of the
original AVaR tail.

